% !TEX program = xelatex
% Biên dịch: xelatex 00-methodology.tex
% Yêu cầu: XeLaTeX + polyglossia + Latin Modern

\documentclass[11pt,a4paper]{article}

%=========================================================
% PACKAGES
%=========================================================
\usepackage[margin=2.2cm]{geometry}
\usepackage{fontspec}
\usepackage{polyglossia}
\setdefaultlanguage{vietnamese}
\setotherlanguage{english}
\setmainfont{Latin Modern Roman}
\setsansfont{Latin Modern Sans}
\setmonofont{Latin Modern Mono}

\usepackage{amsmath}
\usepackage{amssymb}
\usepackage{physics}
\usepackage{xcolor}
\usepackage{graphicx}
\usepackage{tabularx}
\usepackage{booktabs}
\usepackage{longtable}
\usepackage{array}
\usepackage{enumitem}
\usepackage{titlesec}
\usepackage{fancyhdr}
\usepackage{hyperref}
\usepackage{bookmark}
\usepackage{url}
\usepackage{listings}
\usepackage{tcolorbox}
\tcbuselibrary{skins, breakable, listings}

%=========================================================
% STYLE
%=========================================================
\definecolor{accent}{HTML}{0B5FA5}
\definecolor{accent2}{HTML}{B8860B}
\definecolor{danger}{HTML}{B22222}
\definecolor{soft}{HTML}{F4F6FA}
\definecolor{softyellow}{HTML}{FFF8E1}
\definecolor{softred}{HTML}{FDECEC}
\definecolor{codebg}{HTML}{F6F8FA}

\hypersetup{
  colorlinks=true,
  linkcolor=accent,
  urlcolor=accent,
  citecolor=accent2,
  pdftitle={00 - Methodology va Working Protocol - Quantum ML Research},
  pdfauthor={TS. Do Phuc Hao}
}

\titleformat{\section}{\Large\bfseries\color{accent}}{\thesection.}{0.5em}{}
\titleformat{\subsection}{\large\bfseries\color{accent}}{\thesubsection}{0.5em}{}
\titleformat{\subsubsection}{\normalsize\bfseries}{\thesubsubsection}{0.5em}{}

\setlist[itemize]{leftmargin=1.2em, itemsep=2pt, topsep=3pt}
\setlist[enumerate]{leftmargin=1.4em, itemsep=2pt, topsep=3pt}

\pagestyle{fancy}
\fancyhf{}
\rhead{\small\textit{Methodology \& Working Protocol}}
\lhead{\small TS. Đỗ Phúc Hảo}
\rfoot{\thepage}
\renewcommand{\headrulewidth}{0.3pt}

% Hộp ghi chú
\newtcolorbox{note}[1][]{
  enhanced, breakable, colback=soft, colframe=accent,
  boxrule=0.4pt, arc=2pt, left=6pt, right=6pt, top=4pt, bottom=4pt,
  title={#1}, fonttitle=\bfseries\color{white}, coltitle=white,
  colbacktitle=accent
}

\newtcolorbox{rulebox}[1][]{
  enhanced, breakable, colback=softyellow, colframe=accent2,
  boxrule=0.6pt, arc=2pt, left=6pt, right=6pt, top=4pt, bottom=4pt,
  title={#1}, fonttitle=\bfseries\color{white}, coltitle=white,
  colbacktitle=accent2
}

\newtcolorbox{warnbox}[1][]{
  enhanced, breakable, colback=softred, colframe=danger,
  boxrule=0.6pt, arc=2pt, left=6pt, right=6pt, top=4pt, bottom=4pt,
  title={#1}, fonttitle=\bfseries\color{white}, coltitle=white,
  colbacktitle=danger
}

% Listings setup
\lstdefinestyle{codestyle}{
  backgroundcolor=\color{codebg},
  basicstyle=\ttfamily\small,
  breaklines=true,
  frame=single,
  framerule=0.3pt,
  rulecolor=\color{gray!40},
  xleftmargin=0.5em,
  xrightmargin=0.5em,
  aboveskip=4pt,
  belowskip=4pt,
  numbers=none,
  tabsize=2,
  showstringspaces=false
}
\lstset{style=codestyle}

%=========================================================
% METADATA
%=========================================================
\title{\vspace{-1.5cm}\textbf{\color{accent} 00 --- Methodology \& Working Protocol} \\[3pt]
       \Large Hiến pháp dự án nghiên cứu \\[2pt]
       \normalsize \textit{Quantum-Inspired / Quantum Machine Learning} \\[2pt]
       \normalsize \textit{cho Wireless / 6G / NTN Security}}
\author{TS. Đỗ Phúc Hảo \\ \small Khoa CNTT --- Đại học Kiến trúc Đà Nẵng}
\date{Ngày 24 tháng 04 năm 2026 (phiên bản 0.1 --- living document)}

%=========================================================
% DOCUMENT
%=========================================================
\begin{document}
\maketitle
\thispagestyle{fancy}

\begin{note}[Quy ước đọc tài liệu này]
Đây là \textbf{tài liệu ràng buộc} (\textit{binding document}) cho mọi phiên làm việc sau này. Mỗi file phân tích, stress test, deep-dive, hay bản thảo paper đều phải \textbf{tuân theo các quy ước} được đặt ra ở đây. Tài liệu sẽ được cập nhật liên tục (\textit{living document}) khi có thay đổi --- mỗi thay đổi kèm theo bump số hiệu phiên bản ở phụ lục.
\end{note}

\tableofcontents
\vspace{0.4cm}
\hrule
\vspace{0.4cm}

%=========================================================
\section{Mục đích và phạm vi}
%=========================================================

Tài liệu này quy định \textbf{triết lý nghiên cứu, quy trình phân tích, chuẩn trình bày, chuẩn thực nghiệm, và pipeline tính toán} cho dự án chuyển hướng sang Quantum-Inspired / Quantum Machine Learning (QI/QML) cho wireless / 6G / NTN security.

Ba đối tượng sử dụng tài liệu:
\begin{itemize}
  \item \textbf{TS. Đỗ Phúc Hảo} --- tác giả chính, người ra quyết định cuối.
  \item \textbf{Trợ lý nghiên cứu AI} --- mọi phiên làm việc phải tra cứu file này trước khi sinh nội dung.
  \item \textbf{Cộng sự tương lai} --- sinh viên hoặc đồng nghiệp tham gia dự án sau này.
\end{itemize}

Các quy ước được \textbf{extract từ hai ``gems''}:
\begin{itemize}
  \item \textit{Gem 1 --- The Transactions Architect} (triết lý \& output format).
  \item \textit{Gem 2 --- The TikZ Architect} (chuẩn hình vẽ).
\end{itemize}

%=========================================================
\section{Triết lý nghiên cứu (Research Philosophy)}
%=========================================================

\subsection{Paradigm: Optimization $\rightarrow$ Learning}

Mọi bài báo chất lượng IEEE Transactions (TWC, JSAC, T-IFS, IoTJ\dots) mà dự án nhắm đến \textbf{phải} có ba trụ cột sau, \textit{không được thiếu}:

\begin{enumerate}
  \item \textbf{Rigorous Mathematical Foundation.} \\
        Công thức hoá bài toán dưới dạng tối ưu hoá có chứng minh (optimization problem). Trình bày thuật toán cổ điển trước: Block Coordinate Descent (BCD), Majorization-Minimization (MM), Lagrangian Dual, Successive Convex Approximation (SCA), hoặc Alternating Optimization.
  \item \textbf{Realistic Handle on Uncertainty.} \\
        Bất định phải được mô hình hoá tường minh: phân phối xác suất (Gamma, Poisson, Rayleigh, Rician), ước lượng kênh không hoàn hảo ($\hat{\mathbf{h}} = \mathbf{h} + \boldsymbol{\Delta h}$), mobility/Doppler.
  \item \textbf{Justifiable Transition to AI/Quantum.} \\
        Chuyển sang DRL, KAN, QNN, QGNN \textit{chỉ khi} có lý lẽ chính đáng (thường là scale $N > 1000$, real-time constraint, hoặc non-stationary environment). Không bao giờ dùng AI/Quantum vì \textit{``fancy''}.
\end{enumerate}

\subsection{Ba câu hỏi tự vấn trước mọi ý tưởng}

\begin{rulebox}[Checklist trước khi ngồi viết bất kỳ ý tưởng nào]
\begin{enumerate}
  \item Bài toán có thể viết thành optimization $\mathbf{P_1}$ với objective và constraints rõ ràng không? Nếu \textbf{không}: chưa sẵn sàng.
  \item Tôi đã có classical baseline không tầm thường (BCD / SCA / heuristic tốt) trước khi gọi AI/Quantum không? Nếu \textbf{không}: làm nó trước.
  \item Tôi có số liệu / kịch bản cho thấy classical \textit{fail} ở scale lớn hay constraint khắc nghiệt không? Nếu \textbf{không}: AI/Quantum chưa justified.
\end{enumerate}
\end{rulebox}

%=========================================================
\section{The Anti-Toy Filter (Stress Test)}
%=========================================================

Mọi ý tưởng nghiên cứu khi mới đặt ra đều là \textbf{``toy problem''}. Nhiệm vụ đầu tiên là \textit{nâng cấp} nó lên tầm Transactions.

\subsection{Bốn bộ lọc bắt buộc}

\begin{longtable}{@{}p{0.6cm} p{3.5cm} p{10.4cm}@{}}
\toprule
\textbf{\#} & \textbf{Filter} & \textbf{Thao tác áp dụng} \\
\midrule
\endhead
1 & \textbf{Kill Idealism} &
Perfect CSI $\rightarrow$ channel estimation error $\Delta h \sim \mathcal{CN}(0, \sigma^2_e)$;
static nodes $\rightarrow$ high-mobility + Doppler shift;
fixed workload $\rightarrow$ stochastic arrival (Poisson $\lambda$, Gamma service time);
known environment $\rightarrow$ unknown / non-stationary. \\
2 & \textbf{Inject Conflict} &
Bắt buộc có \textit{trade-off không thể né}. Ví dụ: maximize throughput thì phải thêm strict energy budget $E_{\max}$ HOẶC jitter constraint $\sigma_\tau \leq \tau_{\max}$ HOẶC secrecy constraint $R_s \geq R_{s,\min}$. \\
3 & \textbf{Complexity Injection} &
Bài toán phải là \textbf{MINLP} hoặc \textbf{non-convex}: kết hợp biến rời rạc (association, selection, discrete phase shift) với biến liên tục (power, beamforming). Chứng minh NP-hard hoặc ít nhất non-convex. \\
4 & \textbf{Scale Justification} &
Nếu bài gốc nhỏ ($N < 100$), phải mở rộng lên $N > 1000$ (số users, IoT devices, subcarriers, qubits logic, satellite beams\dots) để AI/Quantum là lựa chọn \textit{duy nhất hợp lý}. \\
\bottomrule
\end{longtable}

\subsection{Reviewer-2 lens}

Trước khi submit, tự đóng vai \textbf{Reviewer 2 khắc nghiệt nhất}. Các câu hỏi phải trả lời được:

\begin{itemize}
  \item ``Assumption về Perfect CSI này không đứng được --- đã có paper nào thả assumption này chưa?''
  \item ``Baseline cổ điển của các anh chỉ là random search? Đó là straw-man, reject.''
  \item ``Scale $N = 10$ thì AI/Quantum có gì đặc biệt? Exhaustive search còn nhanh hơn. Reject.''
  \item ``Quantum advantage ở đâu? Số liệu đâu? Reject.''
\end{itemize}

\begin{warnbox}[Cấm kỵ --- sẽ bị reject ngay vòng đầu]
\begin{itemize}
  \item Perfect CSI + không nhắc đến channel estimation.
  \item Single-user scenario mà đi bàn ``massive connectivity''.
  \item Baseline chỉ có random / uniform power allocation.
  \item Quantum circuit 4 qubits trên toy dataset 50 mẫu.
  \item Không có phân tích độ phức tạp / barren plateau / noise resilience.
\end{itemize}
\end{warnbox}

%=========================================================
\section{Output Format cho Problem Formulation}
%=========================================================

Mọi file phân tích niche / stress test sau này phải tuân theo \textbf{cấu trúc 4 mục} sau (lấy từ Gem 1):

\subsection{Mục 1 --- Critique of Original Idea}
\begin{itemize}
  \item Nêu rõ ý tưởng ban đầu (raw idea, 2--3 câu).
  \item \textbf{Reviewer-2 lens}: liệt kê tối thiểu 3 điểm yếu khiến bị reject.
  \item Xác định \textit{toy assumptions} cần loại bỏ.
\end{itemize}

\subsection{Mục 2 --- Upgraded System Model}
\begin{itemize}
  \item Tập hợp (sets): $\mathcal{M}$ (ví dụ: tập users), $\mathcal{N}$ (tập subcarriers), $\mathcal{T}$ (time slots), $\mathcal{K}$ (tập satellites)\dots
  \item Indices: $m, n, t, k$.
  \item Stochastic element: phân phối rõ ràng, tham số bất định.
  \item Channel / traffic / energy / security model: công thức đầy đủ.
\end{itemize}

\subsection{Mục 3 --- The ``Hard'' Problem Formulation}

Trình bày bài toán $\mathbf{P_1}$ bằng LaTeX chuẩn:

\begin{lstlisting}
$$
\begin{aligned}
\mathbf{P_1}: \quad
& \underset{\{x_{m,n}\}, \{p_{m,t}\}}{\text{maximize}}
  && \sum_{m \in \mathcal{M}} \sum_{t \in \mathcal{T}} R_{m,t}(\mathbf{p}, \mathbf{x}) \\
& \text{s.t.}
  && \sum_{t} p_{m,t} \leq E_{\max}, \quad \forall m \in \mathcal{M}, \\
&  && R_{m,t}^{\text{sec}} \geq R_{s,\min}, \quad \forall m, t, \\
&  && x_{m,n} \in \{0, 1\}, \quad \forall m, n, \\
&  && \sum_{m} x_{m,n} \leq 1, \quad \forall n. \quad \text{(coupling)}
\end{aligned}
$$
\end{lstlisting}

Yêu cầu:
\begin{itemize}
  \item Objective \textbf{non-convex hoặc multi-objective}.
  \item Ít nhất \textbf{một coupling constraint} (chọn $x$ giới hạn feasible region của $p$).
  \item Ít nhất \textbf{một biến rời rạc} $\rightarrow$ MINLP.
\end{itemize}

\subsection{Mục 4 --- Dual-Approach Roadmap}

\textbf{Step 1 (The Math)} --- Simplified case (thả 1--2 assumption):
\begin{itemize}
  \item BCD: chia $\mathbf{P_1}$ thành $\mathbf{P_{1a}}$ (fix $x$, tối ưu $p$) và $\mathbf{P_{1b}}$ (fix $p$, tối ưu $x$).
  \item SCA / MM: ghi rõ surrogate function.
  \item Lagrangian Dual: viết dual problem và KKT.
  \item Chứng minh convergence đến stationary point.
\end{itemize}

\textbf{Step 2 (The AI/Quantum)} --- Why classical fails at scale:
\begin{itemize}
  \item Độ phức tạp $\mathcal{O}(N^3)$ hay tệ hơn $\rightarrow$ không chạy real-time.
  \item Non-stationary channel $\rightarrow$ phải re-solve mỗi $T_c$ ms.
  \item Proposed weapon: QGNN / QNN / QDRL / QFL --- nói rõ \textbf{ansatz}, \textbf{encoding}, \textbf{training loss}, \textbf{expected advantage}.
\end{itemize}

%=========================================================
\section{Quy ước LaTeX \& Notation}
%=========================================================

\subsection{Math notation}

\begin{longtable}{@{}p{5cm} p{10cm}@{}}
\toprule
\textbf{Yếu tố} & \textbf{Quy ước} \\
\midrule
\endhead
Inline math & Dùng \texttt{\$...\$}, KHÔNG dùng \texttt{\$\$...\$\$} inline. \\
Display math (1 dòng) & \texttt{\textbackslash[...\textbackslash]}. \\
Display math optimization & \texttt{\$\$\textbackslash begin\{aligned\}...\textbackslash end\{aligned\}\$\$}. \\
Sets (tập hợp) & \texttt{\textbackslash mathcal\{M\}, \textbackslash mathcal\{N\}}. \\
Matrices / vectors & \texttt{\textbackslash mathbf\{H\}, \textbackslash mathbf\{p\}} (bold thẳng). \\
Random variables & \texttt{\textbackslash mathsf\{X\}} hoặc in hoa thường --- thống nhất trong 1 paper. \\
Kỳ vọng / xác suất & \texttt{\textbackslash mathbb\{E\}[\textbackslash cdot], \textbackslash mathbb\{P\}(\textbackslash cdot)}. \\
Complex normal & \texttt{\textbackslash mathcal\{CN\}(0, \textbackslash sigma\textasciicircum 2)}. \\
Quantum states & \texttt{\textbackslash ket\{\textbackslash psi\}, \textbackslash bra\{\textbackslash psi\}, \textbackslash braket\{\textbackslash psi|\textbackslash phi\}} (package \texttt{physics}). \\
Optimization label & \texttt{\textbackslash mathbf\{P\_1\}, \textbackslash mathbf\{P\_\{1a\}\}}. \\
\bottomrule
\end{longtable}

\subsection{Engine LaTeX}

\begin{itemize}
  \item \textbf{File phân tích tiếng Việt} (00, 01, 02\dots): \textbf{XeLaTeX + polyglossia}, font Latin Modern. Header chuẩn: xem file 01.
  \item \textbf{File hình vẽ standalone} (TikZ): \textbf{pdflatex + standalone class} theo Gem 2 (để dễ copy-paste sang Overleaf nếu cần).
  \item \textbf{Paper bản thảo IEEE}: \textbf{pdflatex + IEEEtran.cls}. Không trộn với tiếng Việt.
\end{itemize}

\subsection{Quy ước trích dẫn}

\begin{itemize}
  \item File phân tích: footnote URL là đủ cho phiên khởi đầu; sau này chuyển sang BibTeX.
  \item File paper: BibTeX riêng, format IEEE (\texttt{IEEEtran.bst}).
  \item Tên file bib: \texttt{refs.bib}, dùng key \texttt{<firstauthor><year><keyword>}, ví dụ \texttt{duong2025nonCentralizedQNN}.
\end{itemize}

%=========================================================
\section{Quy ước hình vẽ (Gem 2 --- TikZ Architect)}
%=========================================================

\subsection{Nguyên tắc chung}

\begin{itemize}
  \item Mỗi hình là \textbf{một file \texttt{.tex} standalone riêng} trong thư mục \texttt{figures/}.
  \item Biên dịch bằng \textbf{pdflatex}, xuất PDF, include vào văn bản chính bằng \texttt{\textbackslash includegraphics}.
  \item \textbf{KHÔNG} hard-code toạ độ; dùng \texttt{positioning} (\texttt{below=of NodeA}, \texttt{right=1.5cm of NodeB}).
  \item Tên file: \texttt{fig-<short-slug>.tex} $\rightarrow$ output \texttt{fig-<short-slug>.pdf}.
\end{itemize}

\subsection{Palette chuẩn dự án}

\begin{lstlisting}[language=TeX]
\definecolor{primary}{RGB}{11, 95, 165}     % accent blue
\definecolor{secondary}{RGB}{184, 134, 11}  % accent gold
\definecolor{success}{RGB}{46, 139, 87}     % green
\definecolor{danger}{RGB}{178, 34, 34}      % red
\definecolor{neutral}{RGB}{96, 106, 115}    % gray
\end{lstlisting}

\subsection{Thư viện TikZ bắt buộc load}

\begin{lstlisting}[language=TeX]
\usetikzlibrary{arrows.meta, positioning, calc, shapes,
                fit, backgrounds, shadows.blur, decorations.pathreplacing}
\usepackage{pgfplots}
\pgfplotsset{compat=1.18}
\end{lstlisting}

\subsection{Node styles tiêu chuẩn}

\begin{itemize}
  \item \texttt{process}: hình chữ nhật bo góc, fill trắng, viền xám.
  \item \texttt{decision}: kim cương, fill \texttt{secondary!10}.
  \item \texttt{quantum}: hình chữ nhật đậm viền, fill \texttt{primary!10}, biểu thị khối lượng tử.
  \item \texttt{classical}: tương tự quantum nhưng fill \texttt{success!10}.
  \item \texttt{data}: hình parallelogram, fill \texttt{neutral!10}.
\end{itemize}

\subsection{Danh mục hình dự kiến cho mỗi paper}

\begin{enumerate}
  \item System model (kiến trúc tổng thể).
  \item Pipeline / data flow (end-to-end).
  \item Proposed ansatz / architecture (circuit + NN).
  \item Convergence curve (loss vs epoch / iteration).
  \item Performance comparison (bar chart / line chart).
  \item Ablation heatmap (2D hyperparameter sweep).
  \item Complexity / scalability plot (log-log).
\end{enumerate}

%=========================================================
\section{Pipeline tính toán \& tái sản xuất (Reproducibility)}
%=========================================================

\subsection{Phần cứng dự án}

\begin{longtable}{@{}p{4cm} p{11cm}@{}}
\toprule
\textbf{Nguồn} & \textbf{Chi tiết} \\
\midrule
\endhead
RTX 4090 server & 24 GB VRAM, Ada Lovelace --- máy chạy chính cho training dài, simulation $\leq 30$ qubits (state vector). \\
Google Colab Pro+ & A100 (40 GB khi may mắn), T4/V100 mặc định --- dùng cho hyperparameter sweep và bản trình diễn. \\
IBM Quantum Experience & Free tier để chạy thật trên QPU thực (verify, không training). \\
\bottomrule
\end{longtable}

\subsection{Budget qubit logic (quan trọng)}

Bộ nhớ state vector: $16 \cdot 2^N$ byte (double precision complex).

\begin{longtable}{@{}c c p{10cm}@{}}
\toprule
\textbf{Qubits} & \textbf{VRAM} & \textbf{Ghi chú} \\
\midrule
\endhead
$\leq 20$ & $\leq 16$ MB & Bất cứ đâu, CPU là đủ. \\
24 & 256 MB & Colab free / RTX 4090 mượt. \\
28 & 4 GB & Colab Pro+ / RTX 4090. \\
30 & 16 GB & \textbf{Giới hạn an toàn trên RTX 4090}. \\
32 & 64 GB & Vượt RTX 4090 $\rightarrow$ cần cluster / tensor network approximation. \\
\bottomrule
\end{longtable}

\textbf{Sweet spot thực nghiệm cho paper Q1}: \textbf{12--24 qubits} với mở rộng \textbf{đến 28--30 qubits} ở phần scalability experiment.

\subsection{Software stack cố định}

\begin{lstlisting}[language=bash]
# Core
python==3.11
pytorch==2.3.x  (with CUDA 12.1)
pennylane==0.36.x
pennylane-lightning-gpu==0.36.x   # cuQuantum backend
qiskit==1.1.x
qiskit-aer==0.14.x

# Supporting
numpy==1.26.x
scipy==1.13.x
scikit-learn==1.5.x
matplotlib==3.9.x
seaborn==0.13.x

# Experiment tracking
wandb==0.17.x

# Optim (cho baseline)
cvxpy==1.5.x
gurobipy  # nếu có license
\end{lstlisting}

\subsection{Repo structure chuẩn}

\begin{lstlisting}[language=bash]
quantum-research-<niche>/
|-- README.md
|-- LICENSE
|-- environment.yml       # conda env (pinned)
|-- requirements.txt      # pip fallback
|-- Makefile              # make setup | train | eval | figures
|-- .gitignore            # data/, results/, *.pt, *.ckpt
|-- src/
|   |-- __init__.py
|   |-- models/           # qnn.py, qgnn.py, qfl.py, baselines.py
|   |-- data/             # loaders, transforms
|   |-- training/         # trainer, callbacks, scheduler
|   |-- ansatz/           # VQC ansatz library
|   `-- utils/            # metrics, logging, seeds
|-- experiments/
|   |-- exp01_classical_baseline/
|   |   |-- config.yaml
|   |   `-- run.py
|   `-- exp02_quantum_main/
|-- configs/              # YAML configs gốc
|-- notebooks/
|   `-- colab_bootstrap.ipynb
|-- data/                 # .gitignore, symlink / download script
|-- results/              # .gitignore
|   |-- checkpoints/
|   |-- metrics/
|   `-- figures/
|-- paper/                # LaTeX source
|   |-- main.tex
|   |-- refs.bib
|   `-- figures/
`-- tests/                # pytest
\end{lstlisting}

\subsection{Workflow ``code local $\rightarrow$ GitHub $\rightarrow$ compute''}

\begin{enumerate}
  \item Code \& test trên máy local (CPU / laptop).
  \item Commit + push lên GitHub (repo private).
  \item Trên Colab Pro+: mở \texttt{colab\_bootstrap.ipynb} $\rightarrow$ clone + \texttt{pip install -r requirements.txt} $\rightarrow$ \texttt{python experiments/expXX/run.py --config ...}
  \item Trên RTX 4090 server: \texttt{git pull \&\& conda activate qml \&\& make train EXP=exp02}
  \item W\&B tự log tất cả; checkpoint về \texttt{results/checkpoints/<exp\_name>/<run\_id>/}.
  \item Sau mỗi exp: \texttt{git pull} kết quả về local $\rightarrow$ phân tích $\rightarrow$ vẽ hình $\rightarrow$ viết paper.
\end{enumerate}

\subsection{Reproducibility checklist (bắt buộc cho mỗi exp)}

\begin{itemize}
  \item Random seed cố định, log ra config.
  \item Tối thiểu \textbf{5 seeds khác nhau} cho mỗi số liệu công bố $\rightarrow$ báo cáo $\text{mean} \pm \text{std}$.
  \item Config YAML full, không có hard-code magic number trong code.
  \item Git commit hash log vào W\&B mỗi run.
  \item Dataset version (DVC hoặc hash file).
  \item Training time + GPU model log vào W\&B.
\end{itemize}

%=========================================================
\section{Quy ước đặt tên file}
%=========================================================

\begin{longtable}{@{}p{5.5cm} p{9.5cm}@{}}
\toprule
\textbf{File} & \textbf{Nội dung} \\
\midrule
\endhead
\texttt{00-methodology.tex} & Hiến pháp (file này). \\
\texttt{01-phan-tich-gs-duong-qml.tex} & Phân tích profile GS + 3 niche + survey kiến thức nền. \\
\texttt{02-stress-test-<niche>.tex} & Áp Anti-Toy Filter + 4-part formulation cho niche đã chọn. \\
\texttt{03-sota-<niche>.tex} & Literature deep-dive 15--20 paper, bảng SOTA, gap matrix. \\
\texttt{04-experiment-plan-<niche>.tex} & Dataset, baselines, metrics, ablation, compute budget. \\
\texttt{05-tier-a-reading-note.tex} & Tóm tắt Tier A (4 foundation paper). \\
\texttt{06-tier-b-reading-note.tex} & Tóm tắt Tier B (paper lab GS Duong). \\
\texttt{figures/fig-<slug>.tex} & Hình standalone TikZ, pdflatex. \\
\texttt{paper/draft-v0.1.tex} & Bản thảo paper đầu tiên. \\
\bottomrule
\end{longtable}

%=========================================================
\section{Chuẩn thực nghiệm cho Transactions-grade paper}
%=========================================================

\subsection{Baselines bắt buộc}

Mọi paper QML/QI phải có \textbf{tối thiểu 4 nhóm baseline}:

\begin{enumerate}
  \item \textbf{Classical naive}: random / uniform / greedy --- để thấy non-triviality.
  \item \textbf{Classical optimization mạnh}: BCD / SCA / CVX exact solver cho simplified case.
  \item \textbf{Classical deep learning mạnh}: DNN / GNN / DRL state-of-the-art (cùng kiến trúc tổng thể, chỉ thay phần quantum bằng classical tương đương).
  \item \textbf{Ablations nội bộ}: các biến thể của đề xuất để cô lập đóng góp (xem Mục \ref{subsec:ablation}).
\end{enumerate}

\subsection{Metrics bắt buộc}

\begin{itemize}
  \item \textbf{Performance metric chính}: tùy bài toán (throughput, secrecy rate, F1, accuracy, AUC, loss).
  \item \textbf{Convergence}: loss vs iteration / epoch.
  \item \textbf{Robustness}: performance dưới channel error $\Delta h$ biến thiên, noise level khác nhau.
  \item \textbf{Complexity}: số FLOPs, wall-clock time, số qubit \& circuit depth.
  \item \textbf{Quantum-specific}: expressibility, entanglement capability, barren plateau diagnostic (variance of gradient).
\end{itemize}

\subsection{Ablation study (quan trọng)}
\label{subsec:ablation}

Tối thiểu:
\begin{enumerate}
  \item Số qubits: $\{4, 8, 12, 16, 20\}$.
  \item Circuit depth / số variational layers: $\{2, 4, 6, 8\}$.
  \item Encoding scheme: angle vs amplitude vs re-uploading.
  \item Ansatz family: hardware-efficient vs problem-inspired.
  \item Noise model: noiseless vs depolarizing vs realistic IBM backend.
  \item Training hyperparameters: learning rate, optimizer (Adam, SPSA).
\end{enumerate}

\subsection{Thống kê}

\begin{itemize}
  \item 5 seeds tối thiểu, báo cáo mean $\pm$ std.
  \item Với so sánh: paired t-test hoặc Wilcoxon signed-rank, báo cáo $p$-value.
  \item Confidence interval 95\% trên hình nếu khả thi.
\end{itemize}

%=========================================================
\section{Chuẩn viết bài theo IEEE Transactions}
%=========================================================

\subsection{Cấu trúc paper chuẩn}

\begin{enumerate}
  \item \textbf{Abstract} (200--250 từ): vấn đề, đóng góp (3 bullet), kết quả số.
  \item \textbf{Introduction}: (a) bối cảnh, (b) literature gap rõ ràng, (c) đóng góp 3--5 bullet, (d) tổ chức paper.
  \item \textbf{Related Work}: chia theo nhóm chủ đề, bảng so sánh với các paper gần nhất.
  \item \textbf{System Model}: nêu giả thuyết + notation + channel/traffic model.
  \item \textbf{Problem Formulation}: $\mathbf{P_1}$ với objective + constraints + chứng minh non-convex.
  \item \textbf{Proposed Approach}:
    \begin{itemize}
      \item Classical warm-up (BCD / SCA) cho simplified case.
      \item Proposed QNN / QGNN / QDRL với architecture + encoding + loss + training.
      \item Complexity analysis.
    \end{itemize}
  \item \textbf{Simulation Results}: setup + baselines + main result + ablation + robustness.
  \item \textbf{Conclusion \& Future Work}.
  \item \textbf{References}: IEEE format.
\end{enumerate}

\subsection{Ngôn ngữ \& văn phong}

\begin{itemize}
  \item Viết \textbf{ngắn, chủ động, khẳng định}. Tránh ``it is believed that'', ``one might consider''.
  \item Dùng thì hiện tại cho mô tả phương pháp, quá khứ cho thí nghiệm đã chạy.
  \item Công thức phải được giải thích \textit{bằng lời} trước khi viết ra.
  \item Mỗi hình có caption tự đủ (đọc mình nó hiểu được).
  \item Mỗi bảng có tiêu đề rõ + chú thích dưới bảng.
\end{itemize}

\subsection{Danh sách tạp chí mục tiêu (xếp theo ưu tiên)}

\begin{enumerate}
  \item \textbf{IEEE Transactions on Wireless Communications (TWC)} --- flagship wireless, GS Duong cộng tác thường xuyên.
  \item \textbf{IEEE Journal on Selected Areas in Communications (JSAC)} --- chuyên đề.
  \item \textbf{IEEE Transactions on Information Forensics and Security (T-IFS)} --- hướng security.
  \item \textbf{IEEE Internet of Things Journal (IoTJ)} --- hướng IoT/NTN.
  \item \textbf{IEEE Transactions on Network Science and Engineering (TNSE)} --- GS đang edit.
  \item \textbf{IEEE Open Journal of the Communications Society (OJ-COMS)} --- open access, fast review, là ``cửa vào'' tốt cho paper đầu tiên.
  \item \textbf{IEEE Wireless Communications Letters (WCL)} --- letter 4--5 trang, cho kết quả nhỏ gọn nhưng có điểm.
\end{enumerate}

\textbf{Chiến lược đề xuất}: Paper 1 nhắm \textbf{WCL hoặc OJ-COMS} để có entry point nhanh, paper 2 nhắm \textbf{TWC / JSAC}.

%=========================================================
\section{Nguyên tắc làm việc với trợ lý AI}
%=========================================================

\subsection{Mỗi phiên bắt đầu}
\begin{itemize}
  \item Trợ lý \textbf{phải tra cứu} file này (\texttt{00-methodology.tex}) trước khi sinh nội dung mới.
  \item Nếu có conflict giữa yêu cầu mới và quy ước ở đây: \textbf{hỏi tác giả} trước khi override.
  \item Mọi ý tưởng mới phải đi qua Anti-Toy Filter (Section 3).
\end{itemize}

\subsection{Mỗi output bắt buộc}
\begin{itemize}
  \item Tuân thủ 4-part format (Section 4) nếu là problem formulation.
  \item Tuân thủ LaTeX conventions (Section 5).
  \item Hình vẽ tuân thủ Gem 2 (Section 6).
  \item Trích dẫn đầy đủ khi nêu paper cụ thể.
\end{itemize}

\subsection{Cấm}
\begin{itemize}
  \item Cấm sinh code không chạy / không kiểm thử logic.
  \item Cấm đưa ra công thức mà không giải thích ký hiệu.
  \item Cấm sinh số liệu giả (fake results) --- chỉ placeholder rõ ràng đánh dấu \texttt{[TBD]}.
  \item Cấm dùng emoji trong file LaTeX cuối cùng (dùng \texttt{\textbackslash bigstar} thay thế).
\end{itemize}

%=========================================================
\section*{Phụ lục A --- Lịch sử phiên bản}
\addcontentsline{toc}{section}{Phụ lục A --- Lịch sử phiên bản}
%=========================================================

\begin{longtable}{@{}p{2cm} p{3cm} p{10cm}@{}}
\toprule
\textbf{Version} & \textbf{Ngày} & \textbf{Thay đổi} \\
\midrule
\endhead
0.1 & 2026-04-24 & Phiên bản đầu. Extract từ Gem 1 + Gem 2. Định nghĩa Anti-Toy Filter, 4-part output, LaTeX conventions, compute pipeline, repo structure, reproducibility checklist, writing standards. \\
\bottomrule
\end{longtable}

%=========================================================
\section*{Phụ lục B --- Checklist nhanh trước khi commit một file phân tích}
\addcontentsline{toc}{section}{Phụ lục B --- Checklist nhanh}
%=========================================================

\begin{itemize}
  \item[$\square$] Anti-Toy Filter đã áp đủ 4 bước?
  \item[$\square$] Có ít nhất 1 coupling constraint?
  \item[$\square$] Có biến rời rạc $\rightarrow$ MINLP?
  \item[$\square$] Có classical baseline mạnh ngoài random?
  \item[$\square$] Có lý do cụ thể tại sao cần AI/Quantum (scale, non-stationary, real-time)?
  \item[$\square$] Notation tuân thủ Section 5?
  \item[$\square$] Danh sách hình dự kiến đã liệt kê?
  \item[$\square$] File biên dịch sạch bằng xelatex?
\end{itemize}

\vspace{0.4cm}
\hrule
\vspace{0.3cm}
\textit{\small Phiên bản 0.1 --- 24/04/2026. Mọi phiên bản kế tiếp phải ghi rõ trong Phụ lục A.}

\end{document}
